\Needspace{20\baselineskip}
\section{Relocating specimens and constructing the concrete interface}
\label{sec:concrete-specimen-system}
The context-step field of the abstract specimen interface is now constructed
from the concrete graph semantics. The missing step after synchronized search
was to rebuild an allowed specimen at the successor coordinate. This section
supplies that construction, with recursive fillings and actual endpoint
transitions, and proves the safety and robustness bridges.

Allowed-context membership must respect literal-label graph isomorphism.
Output closure and refocusing of the extrinsic behaviour remain the paper's
semantic hypotheses. No context-step law, successor-specimen existence,
or safety bridge is supplied as a new hypothesis.

There is still a distinction between this interface on named states and the
final contextual theorem. The abstract simulation theorem requires literal
state determinism, whereas the concrete structural theorem proves determinism
modulo recursive isomorphism. A compatible state-class treatment, together
with the concrete observation interface and its initial-lifting theorem,
remains necessary. This section does not replace that obligation by a stronger
determinism assumption on raw names.

\Needspace{14\baselineskip}
\subsection{Exact hole inventories under focus movement}
\begin{lemma}[A focus has no hole occurrence]
Every occurrence address in a focus label observes no hole.
\end{lemma}
\begin{proof}
At every depth the focus label is primitive. Its shallow occurrence is not a
hole and its child-occurrence type is empty. Case analysis on the depth and
address proves the assertion.
\end{proof}

\lean{Graph/FocusContexts}{10}{15}{focus_fiber_no_hole}

\begin{lemma}[Insertion preserves hole occurrences and uniqueness]
Every hole occurrence after focus insertion comes from a retained old edge
with exactly the same label-local address. Consequently insertion preserves
and reflects the existence of every hole and the unique-hole property.
\end{lemma}
\begin{proof}
An occurrence has an outer edge and a label-local address. Split the edge into
an old edge or the fresh singleton focus. The first case retains the address;
the second contradicts the preceding lemma. The converse inclusion embeds old
edges into the left summand. To preserve uniqueness, pull two equal-labelled
hole occurrences back to the old graph, use its uniqueness, and map the equal
addresses forward. Reflection follows from the injectivity of the old-edge map.
\end{proof}

\lean{Graph/FocusContexts}{22}{28}{hole_occurrence}

\lean{Graph/FocusContexts}{31}{37}{hasHole_iff}

\lean{Graph/FocusContexts}{40}{50}{uniqueHoles_iff}

\begin{lemma}[Erasure preserves holes and uniqueness]
Erasing a focussed context preserves every hole in its full recursive inventory
and preserves uniqueness of hole occurrences.
\end{lemma}
\begin{proof}
Retained edges are a subtype of old edges with identical labels. This injective
embedding reflects uniqueness and maps every surviving occurrence back. In the
other direction an old hole cannot lie on the erased focus: transporting its
address along the focus-label equality would contradict the no-hole lemma.
Thus its edge belongs to the retained subtype, with its entire address unchanged.
\end{proof}

\lean{Graph/FocusContexts}{57}{58}{underlying_uniqueHoles}

\lean{Graph/FocusContexts}{61}{71}{underlying_hasHole_iff}

\termdef{iso-closed-contexts}
\begin{definition}[Isomorphism-invariant allowed contexts]
A predicate of allowed erased contexts is isomorphism-invariant if a literal-label
graph isomorphism transports membership from its source to its target.
Applying the law to the inverse also reflects membership. This is invariance
under names, not closure under rewriting.
\end{definition}

\lean{Graph/FocusContexts}{76}{77}{IsoClosedContexts}

\termdef{relocated-context}
\begin{definition}[Context at a new coordinate]
For specimen $S$ and typed coordinate $q$ on its erased context, erase the old
focus and insert the mode of $q$ at its term vertex. The resulting graph is a
unique-hole context, with the newly inserted edge as its distinguished focus.
\end{definition}
\begin{proof}[Verification]
Erasure preserves uniqueness and insertion reflects it back to a focussed
context. The old erasure is focus-free, so insertion supplies the required
unique focus. The insertion--erasure round trip identifies the new erased
context with the old one. Applying allowed-membership invariance to the inverse
of that isomorphism proves the new context allowed. Composing the two inventory
laws proves that exactly the original holes remain, including those in boxes.
\end{proof}

\lean{Graph/FocusContexts}{84}{87}{contextAt}

\lean{Graph/FocusContexts}{90}{92}{focusAt}

\lean{Graph/FocusContexts}{95}{97}{contextAt_underlying_iso}

\lean{Graph/FocusContexts}{100}{102}{contextAt_allowed}

\lean{Graph/FocusContexts}{105}{107}{contextAt_hasHole_iff}

\begin{lemma}[Entering and remoteness after relocation]
The new focus enters a shallow hole precisely when the new coordinate belongs
to that hole's old source list. A remote coordinate therefore gives a remote
focussed context.
\end{lemma}
\begin{proof}
A hole edge cannot be the newly inserted focus. On a retained edge, membership
of the fresh target in the split source list supplies an old source port.
Collapse the equality to identify it with the chosen coordinate. Conversely,
a source occurrence of that coordinate is split to the fresh target. For
remoteness, either the mode is not rewrite, or the coordinate avoids every
hole source; the equivalence just proved excludes entering in the second case.
\end{proof}

\lean{Graph/FocusContexts}{111}{128}{focusAt_entering_iff}

\lean{Graph/FocusContexts}{131}{139}{focusAt_remote}

\Needspace{14\baselineskip}
\subsection{Insertion commutes with simultaneous filling}
\termdef{focus-extended-family}
\begin{definition}[Extending the filling across a focus]
Insert the focus into a focus-free context $C$. Extend a filling family $H$ by
keeping $H(e)$ at every old edge and assigning the new focus its one-edge
corolla. This new corolla has no internal vertices.
\end{definition}

\lean{Graph/FocusExpansion}{12}{12}{inserted}

\lean{Graph/FocusExpansion}{15}{19}{family}

\begin{definition}[Carrier reassociation]
The vertex bijection sends retained context vertices, filling vertices, and the
fresh focus target to the corresponding three components of the graph obtained
by filling first and inserting afterwards. An apparent fourth component,
internal vertices of the focus corolla, is empty. The edge bijection sends old
filling edges to their original expansion pairs and the new corolla's sole
edge to the inserted focus edge.
\end{definition}
\begin{proof}[Bijection laws]
The inverse maps return each component to its original summand. Each inverse
law is a case split. The empty component is eliminated; all indices in the
singleton focus carrier equal zero.
\end{proof}

\lean{Graph/FocusExpansion}{22}{40}{vertices}

\lean{Graph/FocusExpansion}{43}{52}{edges}

\begin{lemma}[Port compatibility]
Retained host ports commute with reassociation, as do split host ports.
A filling source port maps to the split expanded port; a filling target port
maps to the retained expanded port.
\end{lemma}
\begin{proof}
The host identity is a case split on internal, input, and output ports.
For a split host port, test equality with the chosen coordinate. Host
injectivity preserves this test through expansion. For a filling source,
an internal port cannot equal a retained host port, an input is attached to
the corresponding split host source, and an output is excluded by source
well-formedness. The target proof is dual: an internal port is retained, an
output is attached to the old target, and an input is impossible. The proof
uses only incidence-side injectivity and therefore does not exclude self-looping
attachments by assuming a globally injective attachment map.
\end{proof}

\lean{Graph/FocusExpansion}{55}{58}{host}

\lean{Graph/FocusExpansion}{61}{70}{split}

\lean{Graph/FocusExpansion}{73}{91}{source_port}

\lean{Graph/FocusExpansion}{94}{107}{target_port}

\begin{theorem}[Commutation of focus insertion and filling]
Expanding the inserted context by the extended family is isomorphic to
inserting the same focus at the corresponding host coordinate of $C[H]$.
\end{theorem}
\begin{proof}
Use the two carrier bijections. Internal and edge labels agree on each
summand. For an old filling edge, map its source and target lists twice and
apply the corresponding port identities pointwise. For the fresh corolla,
its source and target lists are singletons; they map to the retained chosen
port and the fresh target respectively. These are precisely the inserted
focus's incidence lists. Thus all fields of a graph isomorphism are proved.
\end{proof}

\lean{Graph/FocusExpansion}{110}{158}{iso}

\Needspace{14\baselineskip}
\subsection{The successor specimen and full-machine traversal}
\begin{lemma}[Extending recursive certificates]
Every old retained, filled-hole, or boxed filling certificate extends to the
inserted context. The new focus is a retained corolla on both sides. If the
old expansion is hole-free, so is the new expansion.
\end{lemma}
\begin{proof}
Inspect an old edge certificate. Retained copies keep the same edge, label,
and incidence proofs; filled holes keep their related replacements; boxed
edges keep their recursive body-filling certificate and rebuilt box copies.
The new edge has the explicit singleton-corolla certificate. Closure follows
by transporting insertion's hole-free result through the commuting isomorphism.
\end{proof}

\lean{Graph/SpecimenRelocation}{14}{21}{edge_lift}

\lean{Graph/SpecimenRelocation}{24}{31}{edges_lift}

\lean{Graph/SpecimenRelocation}{34}{37}{closed}

\termdef{relocated-specimen}
\begin{definition}[Relocated specimen]
The specimen at $q$ consists of the newly focussed context, its proved allowed
membership, and the two extended families obtained from the old erased
families. The extended certificates and closure proofs supply all remaining
specimen fields. No new relation pair is introduced.
\end{definition}
\begin{lemma}[Endpoint identification and remoteness]
Its two actual endpoints are respectively isomorphic, preserving focus, to the
previously constructed left and right coordinate realizations. If $q$ is
remote, so is the relocated specimen.
\end{lemma}
\begin{proof}
Use insertion--expansion commutation on each side. A graph isomorphism between
uniquely focussed graphs preserves the focus by its label and uniqueness,
yielding the focussed isomorphisms. The remote assertion is the earlier
coordinate-to-context result applied to the constructed focus.
\end{proof}

\lean{Graph/SpecimenRelocation}{46}{57}{relocate}

\lean{Graph/SpecimenRelocation}{60}{63}{relocate_left_iso}

\lean{Graph/SpecimenRelocation}{66}{69}{relocate_right_iso}

\lean{Graph/SpecimenRelocation}{72}{75}{relocate_remote}

\begin{theorem}[Allowed successor specimen for search]
For an output-closed template, a settled specimen with jointly rooted endpoints
and a search successor on either side admits a remote allowed specimen reached
by search on both sides. Recursive renaming of the originally supplied target
is permitted.
\end{theorem}
\begin{proof}
The checked traversal theorem supplies a synchronized typed coordinate and
proves it remote. Relocate the specimen there. Compose each search endpoint
with the inverse of its relocation isomorphism to obtain the two transitions
into the new specimen.
\end{proof}

\lean{Graph/SpecimenRelocation}{79}{93}{settled_successor}

\begin{lemma}[Settled full-machine steps are search]
A full-UAM step from a settled specimen's endpoint must be search.
\end{lemma}
\begin{proof}
Split the transition into search or rewrite. The latter requires a rewrite
source, contradicting either backtrack mode or search mode in the settled
hypothesis. No locality or determinism assumption on the behaviour is used.
\end{proof}
\begin{theorem}[Full-UAM context-step law]
Under allowed-membership invariance and output closure, a jointly rooted
settled specimen with a full-UAM successor on either side has a remote allowed
successor specimen reached by full-UAM steps on both sides.
\end{theorem}
\begin{proof}
Joint rootedness and output closure exclude the exiting boundary; settledness
excludes the entering search boundary. Reduce whichever full-machine step is
available to search, reflect it to the erased context, and obtain a coordinate
$q$. Its search-step theorem proves it remote. Build the relocated specimen.
The literal search-reflection equivalences, composed with the two literal
relocation isomorphisms, give actual search steps into its endpoints. Inject
these into the full-UAM relation. The conclusion uses actual named states,
not a substituted assumption about reachability or a quotient-only step.
\end{proof}

\lean{Graph/SpecimenRelocation}{96}{104}{settled_step_search}

\lean{Graph/SpecimenRelocation}{108}{136}{full_context_step}

\begin{theorem}[Constructed initial specimens]
Relocating any allowed specimen to the unique input produces a specimen with
two initial, hence rooted, endpoints and search mode. The old specimen need
not have rooted endpoints.
\end{theorem}
\begin{proof}
The inserted focus has search mode and an input source. The expansion boundary
law transfers initiality to the left endpoint, and the specimen boundary
agreement transfers it to the right. Zero-step initial runs establish both
rootedness claims. Search mode establishes remoteness. This constructs an
initial witness from an existing specimen; it does not yet construct the
observation/context witnesses for arbitrary related fragments.
\end{proof}

\lean{Graph/SpecimenRelocation}{141}{148}{relocate_initial}

\lean{Graph/SpecimenRelocation}{152}{158}{initial_witness}

\Needspace{14\baselineskip}
\subsection{The concrete abstract-interface instance and its bridges}
\termdef{concrete-specimen-system}
\begin{definition}[Constructed specimen system]
For behaviour $B$, use actual recursive specimens as the abstract specimen
type and their actual filled states as endpoints. Use the context's focus mode
and entering predicate. Rootedness is the existing concrete predicate. Supply
rooted-step preservation from the full-UAM structural theorem, final agreement
from the boundary theorem, and context traversal from the newly proved
full-machine context-step theorem.
\end{definition}
The parameters are an isomorphism-invariant allowed-context class,
output closure of the template at the chosen depth, and refocusing of the
extrinsic behaviour. All operational interface fields are now derived.

\lean{Graph/ConcreteSpecimenSystem}{16}{25}{system}

\begin{lemma}[Quasi-specimen bridge]
A concrete non-rewrite quasi-specimen supplies the abstract rooted
quasi-specimen predicate for the constructed system.
\end{lemma}
\begin{proof}
After receiving both endpoint-rootedness hypotheses, apply the existing
rooted-witness theorem. Its witness has non-rewrite mode, rooted filled
endpoints, and exactly the two required comparison relations. The concrete
unconditional witness and mode agreement have not been discarded from the
original definition; this theorem derives the consequence used by the interface.
\end{proof}

\lean{Graph/ConcreteSpecimenSystem}{29}{33}{quasi_bridge}

\begin{theorem}[Input-safety bridge]
The unrestricted concrete input-safety predicate implies the abstract
input-safety law of the constructed specimen system.
\end{theorem}
\begin{proof}
Apply concrete safety at the entering search specimen. Its stuck alternative
is already the first abstract alternative. Its direct alternative has the
same specimen, settled predicate, exact run lengths, and cost comparison.
For the third alternative retain both runs and their cost comparison and
apply the quasi-specimen bridge. In particular the equivalence relations in
this alternative are preserved, not weakened prematurely to refinement.
\end{proof}

\lean{Graph/ConcreteSpecimenSystem}{36}{44}{inputSafe_bridge}

\begin{theorem}[Robustness bridge]
Concrete robustness implies the abstract robustness law of the constructed
system.
\end{theorem}
\begin{proof}
The abstract source has rewrite mode, so its full-machine transition cannot
be search and must be a rewrite. Apply the concrete rooted robustness theorem:
the unrooted alternative contradicts the two given rootedness hypotheses.
The stuck-tail case is already the abstract case. In the matching case keep
the tail lengths, right run, and cost comparison and apply the quasi-specimen
bridge with the supplied refinement relations.
\end{proof}
\begin{lemma}[Exact concrete lifting]
The abstract lifting relation for this system is precisely existence of a
remote concrete specimen with rooted endpoints equal to the given states.
\end{lemma}
\begin{proof}
Unfold the constructed system and the two definitions. The predicates coincide
definitionally.
\end{proof}

\lean{Graph/ConcreteSpecimenSystem}{48}{59}{robust_bridge}

\lean{Graph/ConcreteSpecimenSystem}{62}{66}{lifting_iff}

\Needspace{14\baselineskip}
\subsection{Nested-hole regression and remaining obligations}
Use the earlier specimen whose sole hole lies strictly inside a box and whose
outer focus searches at the consumer. All contexts are allowed in this
regression, making membership invariance immediate. Relocate it to the
consumer's actual backtrack successor. The inventory theorem retains the
nested hole, the context carries a uniqueness proof, and both filled endpoints
are closed. Applying the two literal search equivalences to the known consumer
step gives actual transitions into the rebuilt specimen on both sides.
Its coordinate is remote, so the rebuilt specimen is remote.

\lean{Graph/SpecimenRelocationExamples}{8}{8}{allowed}

\lean{Graph/SpecimenRelocationExamples}{10}{10}{moved}

\lean{Graph/SpecimenRelocationExamples}{13}{14}{nested_hole}

\lean{Graph/SpecimenRelocationExamples}{17}{17}{unique}

\lean{Graph/SpecimenRelocationExamples}{20}{21}{closed}

\lean{Graph/SpecimenRelocationExamples}{24}{30}{paired_steps}

\lean{Graph/SpecimenRelocationExamples}{33}{34}{remote}

The same example was previously proved unrooted. Moving its focus to the input
now supplies a genuine initial, rooted specimen, exercising construction
without assuming the old source rooted. Finally, instantiate the concrete
system with the empty behaviour: its refocusing condition is vacuous, output
closure is the earlier template theorem, and the nested specimen witnesses
that the constructed abstract specimen type is nonempty.

\lean{Graph/SpecimenRelocationExamples}{38}{42}{initial_rooted}

\lean{Graph/SpecimenRelocationExamples}{45}{46}{concreteSystem}

\lean{Graph/SpecimenRelocationExamples}{48}{48}{interface_nonempty}

\paragraph{What this milestone closes.}
The synchronized-coordinate packaging, unique-hole preservation,
isomorphism-invariant membership transport, full-UAM context-step field, and
concrete safety and robustness bridges are complete. Appendix~\ref{app:remaining-work}
retains only the unfinished portions. Initial observation lifting still needs
its general context and endpoint witnesses. The named-state interface must
also be connected to a deterministic state-class machine before applying the
abstract theorem under the paper's original determinism hypothesis. The
planning estimate is now about 91\% of mathematical effort (plausible range
78--95\%); it remains a judgment rather than a theorem or declaration count.
